Florence vient au lycée avec un sac à dos contenant son matériel scolaire. Sa
masse étant égale à $m=\qty{8.0}{\kg}$, calculer le travail du poids $\vec{P}$
du sac à dos ainsi chargé lorsque~:
\begin{enumerate}
    \item Florence parcourt un couloir rectiligne et horizontal de longueur
          $L=\qty{50}{\m}$;
    \item Florence, initialement au rez-de-chaussée, monte au troisième étage.
\end{enumerate}
\begin{donnees}
    \begin{itemize}
        \item intensité de la pesanteur~: $g=\qty{10}{\N\per\kg}$
        \item hauteur d'un étage~: $H=\qty{2.80}{\m}$.
    \end{itemize}
\end{donnees}

